\documentclass[11pt]{article}
\usepackage[margin=1.1in]{geometry}
\usepackage{amsmath,amssymb}
\usepackage{booktabs}
\usepackage[T1]{fontenc}
\usepackage{lmodern}
\usepackage{microtype}
\usepackage[hidelinks]{hyperref}
\usepackage{url}

\title{\bf Identifiability and Predictability Are Distinct:\\
Measured Preimages Across Generative Systems}
\author{Greg Ward\\\small\texttt{greg@smartledger.technology}}
\date{September 2026}

\newcommand{\bits}{\ensuremath{\log_2|P(O)|}}
\newcommand{\HH}{\ensuremath{\mathcal{H}}}

\begin{document}
\maketitle

\begin{abstract}
A recurring proposal in complexity science is to search for the smallest rule that
generates an observed structure. We argue that this question is not falsifiable as
posed, and replace it with one that is: for a hypothesis space $\HH$ of candidate
generators and an observation $O$, how large is the preimage
$P(O)=\{h\in\HH : O(h)=O(g)\}$, and how does it shrink?
We first measure this exactly---by enumeration, not sampling---in three deterministic
systems: expressions over $\{1,+,\times\}$, the 256 elementary cellular automata, and a
129-rule family of hypergraph rewriting systems; we then test the resulting framework in
a pre-registered fourth system of a deliberately dissimilar kind.

Three results follow. First, the unresolved information decomposes as
\emph{observation deficiency} plus \emph{degeneracy of $\HH$}; the first term goes to
zero in two of our systems and the second is $1.09$ bits of irreducible residual in
the third. Second, the deficiency term is fixable only where the observer may
intervene: given the same automata with the initial state chosen for us rather than
by us, a late-arriving observer is left with $2.96$ bits and identifies the rule only
$20.6\%$ of the time, against $0.00$ bits and a guarantee for an observer permitted to
set the state. Third, and most consequentially, \emph{identifiability and
predictability are distinct, non-redundant properties}: neither determines the other. Elementary cellular automata are identifiable in a
single linear pass and yet unpredictable, their generator-to-behaviour map admitting
no local structure; hypergraph rewriting under a fixed scheduler is predictable
without being identifiable. These examples populate off-diagonal regimes that are
not usually examined together within a common finite-preimage framework, separating
generator recovery from behavioural prediction.

A pre-registered replication on a fourth and deliberately dissimilar system---all
$1{,}500{,}625$ partially observed $4$-state Markov chains under a balanced lumping---meets
all five frozen criteria. There the same separation appears for an entirely different
reason: the generator is exactly known and the outcome still is not, because the law is
stochastic rather than because its behaviour is undecidable. In that system not one generator
is uniquely identified by the pre-registered passive observation map at horizon $T=6$,
while $82.0\%$ of them change quadrant once the observer is allowed to set the initial
state.

We further find that the forward map is not always a function: identical primitives,
rules and iteration counts produce provably non-isomorphic outcomes under different
match orders in $2$ of $8$ rewriting systems tested, so any inverse problem must
carry the scheduler as an explicit argument.
\end{abstract}

\section{Introduction}

That simple rules generate complex behaviour is not in dispute and has not been for
decades. The proposal that recurs, and that motivates this work, is the converse: that
one might systematically recover the simple rule beneath a complex object, and that a
research programme could be organised around doing so.

We think that programme fails as stated, for a reason that is easy to demonstrate and
easy to overlook. Consider the smallest non-trivial generative system available:
expressions built from the literal $1$ under $+$ and $\times$. The cost of the integer
$107$---the minimum number of $1$s needed to write it---is $16$. There is no known
closed form for this cost function; whether $\mathrm{cost}(2^k)=2k$ for all $k$ is
open; the cost is not monotone, falling $49$ times in $[1,200]$ as the target
increases, with $\mathrm{cost}(107)=16$ exceeding $\mathrm{cost}(128)=14$. And the
minimum is not unique: $107$ admits $26{,}624$ minimal expression trees. If the
simplest generative system in existence has no compact, unique, or even stable inverse,
``find the simple rule'' cannot be the research question.

The question that survives is about \emph{preimages}. Fix a hypothesis space $\HH$ of
candidate generators and an observation map $O$. The object of study is
\[
  P(O) \;=\; \{\,h \in \HH \;:\; O(h) = O(g)\,\},
\]
the set of generators indistinguishable from the true one under $O$, and the quantity
of interest is $\bits$, the bits of generator identity still unresolved.
Identifiability is $|P(O)|=1$.

This framing is not new---it is structural identifiability in the sense of Bellman and
{\AA}str\"om~\cite{bellman1970}, a Markov equivalence class in the sense of Verma and
Pearl~\cite{verma1990}, and the Myhill--Nerode construction in the sense of
Nerode~\cite{nerode1958}. Our contribution is not the preimage framing itself, but exact
measurements across several generative regimes, followed by a pre-registered test in a
deliberately dissimilar stochastic system.

\paragraph{Contributions.}
\begin{enumerate}
\item A measured decomposition of unresolved information into a term removable by
better experiment and a term that is not (\S\ref{sec:decomp}), with both terms
computed exactly rather than estimated.
\item A demonstration that the removable term is removable \emph{only under
intervention}, with the passive case measured exactly over every initial condition
(\S\ref{sec:passive}).
\item Evidence that identifiability and predictability are distinct, non-redundant
axes, with all four quadrants populated by measured systems (\S\ref{sec:quadrants}).
\item A pre-registered replication on a fourth system---a family of $1{,}500{,}625$
partially observed Markov chains---that meets all five frozen criteria
(\S\ref{sec:stochastic}), and in which the unpredictability is aleatoric rather than
computational.
\item A negative result of practical consequence: the generator-to-behaviour map for
elementary cellular automata has no local structure, so search in rule space cannot be
guided by proximity (\S\ref{sec:landscape}).
\end{enumerate}

\section{Framework}\label{sec:framework}

Let $\HH$ be a finite set of candidate generators, $G$ a random element of $\HH$ under
a uniform prior, and $O$ an observation map. Writing $n_c = |O^{-1}(c)|$ for the size of
the class $c \in \mathrm{im}\,O$, every generator in a class is equiprobable given that
class, so the residual uncertainty about generator identity is the conditional entropy
\begin{equation}\label{eq:hgo}
  H(G \mid O) \;=\; \frac{1}{|\HH|}\sum_{c \in \mathrm{im}\,O} n_c \log_2 n_c
  \;=\; \mathbb{E}\!\left[\log_2 |P(O)|\right].
\end{equation}
This is $0$ exactly when $O$ separates every generator, and $\log_2|\HH|$ when $O$ is
constant. Every figure reported below as ``bits unresolved'' is this quantity.

We flag one trap explicitly, having fallen into it ourselves. The superficially similar
$\log_2 \mathbb{E}[|P(O)|] = \log_2\big(|\HH|^{-1}\sum_c n_c^2\big)$ is \emph{not}
equal to \eqref{eq:hgo}; by Jensen's inequality it is an upper bound, and the two
differ by up to half a bit in our data. It is the wrong quantity here because it is the
log of an expected \emph{count}, not an expected \emph{code length}, and does not
compose additively the way the decomposition below requires.

\paragraph{Per-generator coordinates.}
Equation~\eqref{eq:hgo} is a single scalar for a whole experiment---an average over the
random generator. The claims of this paper are about how quantities \emph{vary between
generators}, so they need per-generator coordinates, and we keep the two levels
notationally distinct. For an individual $g \in \HH$ write
\[
  i_O(g) \;=\; \log_2\big|P_O(g)\big|,
  \qquad\text{so that}\qquad
  H(G \mid O) \;=\; \mathbb{E}_G\big[\,i_O(G)\,\big],
\]
the mean unresolved generator information being simply the population average of the
per-generator coordinate. Recovery is one question and prediction is another, so alongside
it write
\[
  p_k(g) \;=\; H\big(Y_{t+k} \,\big|\, O_{\le t},\, G = g\big),
\]
where $Y$ is a declared predictive target and $k$ a declared horizon. Then $i_O$ is
epistemic uncertainty about which mechanism is acting, and $p_k$ is uncertainty about
what that mechanism will do \emph{with the mechanism known}: different conditional
entropies answering different questions.

Both coordinates are relative to declarations, and the observation horizon is one of
them. Everything here is stated with respect to a hypothesis space $\HH$, an observation
map $O$---including how much of the process $O$ sees---and an admissible class
$\mathcal{A}$; none of the three may be left implicit. In particular, for a stochastic
process observed to horizon $T$, equality of the finite-dimensional law
$\mathcal{L}(Y_{0:T-1} \mid g_1) = \mathcal{L}(Y_{0:T-1} \mid g_2)$ does not imply equality
of the full process law, so a preimage measured at finite $T$ is an upper bound on the
preimage under unlimited passive data and must be reported with $T$ attached.

Every two-axis statement below is about the pair $\big(i_O(g), p_k(g)\big)$, not about
the population scalars, and the thesis is the pair of non-implications
\begin{equation}\label{eq:nonimp}
  i_O(g_1) = i_O(g_2) \;\;\not\Longrightarrow\;\; p_k(g_1) = p_k(g_2),
  \qquad
  p_k(g_1) = p_k(g_2) \;\;\not\Longrightarrow\;\; i_O(g_1) = i_O(g_2),
\end{equation}
which is what criteria C2 and C3 of \S\ref{sec:stochastic} test directly. Where the
predictive target is deterministic, as for polynomials, $p_k \equiv 0$ and only $i_O$ has
content; \S\ref{sec:stochastic} works in a setting where both vary.

\paragraph{The decomposition, made exact.}
Let $\mathcal{A}$ be the family of observations the experimenter is actually permitted
to perform---the \emph{admissible class}, which encodes what may be intervened on---and
let
\[
  O^\star \;=\; \operatorname*{arg\,min}_{O' \in \mathcal{A}} H(G \mid O') .
\]
Define the \emph{intrinsic degeneracy under $\mathcal{A}$} and the
\emph{observation deficiency of $O$} as
\[
  D_\HH(\mathcal{A}) \;=\; H(G \mid O^\star),
  \qquad
  \Delta_O \;=\; H(G \mid O) - H(G \mid O^\star) .
\]
Then trivially, and with both terms non-negative by construction,
\begin{equation}\label{eq:decomp}
  H(G \mid O) \;=\; \underbrace{\Delta_O}_{\text{a bad experiment}}
  \;+\; \underbrace{D_\HH(\mathcal{A})}_{\text{distinct generators, identical objects}} .
\end{equation}
The content of \eqref{eq:decomp} is not the identity, which is definitional, but the
measured values of its two terms and---critically---the fact that
$D_\HH(\mathcal{A})$ moves when $\mathcal{A}$ does. Enlarging the admissible class from
passive observation to intervention can convert intrinsic degeneracy into mere
deficiency; \S\ref{sec:passive} measures that conversion.

\section{Systems and method}\label{sec:systems}

All results are exact enumerations over finite hypothesis spaces. Nothing below is
sampled unless explicitly stated.

\paragraph{S1: expressions over $\{1,+,\times\}$.}
Leaves are the literal $1$; the budget is the number of leaves. We enumerate every
value reachable at each budget and record, for each integer, the minimum budget at
which it appears.

\paragraph{S2: elementary cellular automata.}
A rule is $8$ bits---one output per three-cell neighbourhood---so $|\HH|=256$ and there
are exactly $8$ bits to recover. Rules act on a periodic ring; we use ring size $64$ for
observation-length curves and $12$ (giving $4096$ states) where exhaustive state-space
analysis is required. A rule is constrained by an observation only where that
observation exercises its neighbourhoods; unexercised entries remain free, so the
preimage is exactly $2^{8-c}$ for coverage $c$.

\paragraph{S3: hypergraph rewriting.}
States are sets of ordered tuples; a rule rewrites a matched sub-hypergraph. One step
applies a maximal non-overlapping set of matches, selected in an order determined by a
scheduler seed. Graph comparison uses a Weisfeiler--Leman
hash~\cite{weisfeiler1968}, a sound but incomplete invariant: a hash difference
proves non-isomorphism, while equality is strong evidence rather than proof. All
residual figures for S3 are therefore \emph{lower bounds} on distinguishability.

\paragraph{Validation.}
Three independent checks against known results. Integer complexity in S1 reproduces
OEIS A005245~\cite{oeis005245} exactly for $n=1,\dots,32$. In S2 the reversible rules
are recovered as $\{15,51,85,170,204,240\}$, and the $256$ rule numbers collapse to
$88$ equivalence classes under reflection and colour swap. In the Boolean-closure
experiment, $\{\wedge,\vee\}$ generates exactly $18$ of the $256$ functions of three
variables, the free distributive lattice on three generators.

\section{Generation is cheap and already classified}\label{sec:generation}

We record these results mainly to set them aside.

Growth class changes discontinuously when a single primitive is added. Over $18$
leaves, $\{1,+\}$ reaches exactly $n$ values at budget $n$; adding $\times$ takes the
count to $404$ and the per-step growth ratio from $1.06$ to $1.37$. For polynomials,
$\{x,1,+\}$ reaches only degree $\le 1$ at any budget, while $\{x,1,+,\times\}$ reaches
$12{,}120$ polynomials at budget $10$, degrees $0$ through $5$.

The sharpest transition is Boolean. On three variables, $\{\wedge,\vee\}$ reaches $18$
of $256$ functions and is stuck there---the monotone functions. $\{\oplus,\wedge\}$
reaches $128$, the $\mathbb{F}_2$ polynomials with zero constant term. Adding the
single constant $1$, with no new operation, reaches all $256$: universality crossed by
one bit.

\begin{table}[t]
\centering
\begin{tabular}{lrr}
\toprule
Gate set & Functions reached (of 256) & \\
\midrule
$\{\wedge\}$ & 7 & \\
$\{\oplus\}$ & 8 & \\
$\{\wedge,\vee\}$ & 18 & monotone functions \\
$\{\oplus,\wedge\}$ & 128 & zero constant term \\
$\{\oplus,\wedge,1\}$ & 256 & universal \\
$\{\wedge,\neg\}$ & 256 & universal \\
$\{\uparrow\}$ & 256 & universal \\
\bottomrule
\end{tabular}
\caption{Boolean closure on three variables. The step from $128$ to $256$ costs one
constant and no new operation.}
\end{table}

None of this is new. Post's lattice~\cite{post1941} classifies every Boolean clone
completely, and has since 1941. We report these figures only to establish that the
apparatus is calibrated, and to make the point that a programme organised around
``which structures does a primitive set generate'' is re-entering a field that has been
occupied for eighty-five years.

\section{The identifiability curve and its plateau}\label{sec:decomp}

\subsection{Cellular automata}

Figure-free summary: we compute the mean over all $256$ rules of $\bits$ as a function
of observation length, for five initial conditions on a $64$-cell ring.

\begin{table}[h]
\centering
\begin{tabular}{lrrrrrrrr}
\toprule
Initial condition & $T{=}1$ & $2$ & $3$ & $5$ & $10$ & $25$ & $100$ & $400$ \\
\midrule
single cell   & 4.00 & 1.88 & 1.66 & 1.60 & 1.60 & 1.60 & 1.58 & 1.58 \\
density $0.02$ & 4.00 & 1.88 & 1.53 & 1.45 & 1.45 & 1.45 & 1.45 & 1.45 \\
density $0.10$ & 4.00 & 1.88 & 1.53 & 1.45 & 1.44 & 1.44 & 1.44 & 1.44 \\
density $0.25$ & 3.00 & 0.84 & 0.81 & 0.81 & 0.81 & 0.81 & 0.81 & 0.81 \\
density $0.50$ & 0.00 & 0.00 & 0.00 & 0.00 & 0.00 & 0.00 & 0.00 & 0.00 \\
\bottomrule
\end{tabular}
\caption{Mean unresolved bits (of $8$) against observation length. The curves flatten
by $T=3$; the height at which they flatten is set by the initial condition.}
\label{tab:plateau}
\end{table}

Two readings. A random half-dense row identifies every rule exactly after a
\emph{single} step. A single-cell seed leaves $1.58$ bits after $400$ steps, having
left $1.66$ after three: running $130$ times longer buys $0.08$ bits. On this ring,
$132$ of $256$ rules are permanently unidentifiable from a single-cell seed at any
observation length. (The count is mildly width-sensitive: a $201$-cell ring gives
$134$, so it should be reported with the ring size attached.)

One well-designed step beats four hundred badly designed ones. This is a demonstration
of infomax experimental design in the sense of Lindley~\cite{lindley1956}, not a new
principle; what is useful here is that the effect is measured exactly rather than
argued.

\subsection{Polynomials: the ideal curve and its failure modes}

Let $\HH$ be integer polynomials of degree $\le d$ with coefficients in $[-C,C]$, and
let $O$ be evaluation at $k$ points. For $d=4$, $C=2$ ($|\HH|=3125$, $11.61$ bits),
consecutive points give $11.61 \to 9.29 \to 5.75 \to 1.06 \to 0$. Degenerate designs
never descend: sampling $x=1$ repeatedly is pinned at $7.90$ bits forever, and two
points repeated at $5.75$.

The threshold arrives earlier than Lagrange interpolation would suggest, and the reason
is instructive. Two candidates agree at $k$ points only if their difference---whose
coefficients lie in $[-2C,2C]$---has $k$ roots there, and a small coefficient box
cannot supply the large coefficients such a product requires.

\begin{table}[h]
\centering
\begin{tabular}{rrrrr}
\toprule
Degree & $C$ & $|\HH|$ & Points to identify & Lagrange bound \\
\midrule
3 & 1 & 81 & 3 & 4 \\
3 & 2 & 625 & 4 & 4 \\
4 & 1 & 243 & 4 & 5 \\
4 & 4 & 59{,}049 & 4 & 5 \\
6 & 1 & 2{,}187 & 4 & 7 \\
6 & 2 & 78{,}125 & 4 & 7 \\
\bottomrule
\end{tabular}
\caption{Shrinking the hypothesis class buys data points. Identifiability is a joint
property of $\HH$ and $O$, never of the object alone.}
\end{table}

Taking this further, with the design chosen \emph{optimally} rather than merely
consecutively, a single point suffices: for $|x| \ge 2C+1$ evaluation is injective,
being positional notation in base $x$ with digits in $[-C,C]$. Thus one sample at
$x=5$ recovers all $11.61$ bits, while $x=1$ leaves $7.90$ bits unresolved at any
multiplicity. A single well-chosen observation beats infinitely many badly chosen ones,
and this is a theorem rather than a tendency.

\subsection{Rewriting: a residual that does not vanish}

$\HH$ is $129$ rewrite rules with left-hand side $\{xy\}$ and a right-hand side of one
to three edges over $x,y,z$ ($7.01$ bits). Each ran five steps from a triangle seed
under a fixed scheduler seed.

\begin{table}[h]
\centering
\begin{tabular}{lrr}
\toprule
Observation & Classes & Bits unresolved \\
\midrule
final $|E|$ only & 11 & 4.02 \\
$|E|$ trajectory & 11 & 4.02 \\
$|E|$ and $|V|$ trajectory & 13 & 3.87 \\
final degree sequence & 35 & 2.16 \\
final graph (WL hash) & 78 & 1.09 \\
full graph trajectory & 78 & 1.09 \\
\bottomrule
\end{tabular}
\caption{Watching the whole trajectory adds nothing over the final state, twice over.
The curve bottoms out at $1.09$ bits.}
\end{table}

\subsection{The decomposition}

\begin{table}[h]
\centering
\begin{tabular}{lrrr}
\toprule
System & $|\HH|$ (bits) & Best observation & Residual \\
\midrule
Elementary cellular automata & 8.00 & 0.00 & 0.00 \\
Polynomials, $d\le4$, $C=2$ & 11.61 & 0.00 & 0.00 \\
Hypergraph rewriting & 7.01 & 1.09 & 1.09 \\
\bottomrule
\end{tabular}
\caption{Equation~\eqref{eq:decomp} evaluated. For the first two systems every
unresolved bit was a bad experiment. For the third, $1.09$ bits survive perfect
observation: distinct rules generating identical objects.}
\end{table}

\section{The forward map is not always a function}\label{sec:scheduler}

A formulation of the generative hypothesis we have seen proposed is $C=G(P,R,I)$:
complexity as a function of primitives, relations, and iteration. In S3 this is false.

Applying the same rule to the same seed for the same number of steps, varying only the
order in which matches are applied, produced provably non-isomorphic outputs in $2$ of
$8$ rules tested---five distinct outcomes in twelve runs for one rule, three for
another. The forward map therefore requires the scheduler as an argument,
\[
  C = G(P,R,I,S),
\]
and whether a given rule is order-independent---causal invariance, in the sense
of~\cite{wolfram2020}---is not decidable in general, confluence being undecidable even for
restricted rewriting classes~\cite{jacquemard2003}, though decidable for others. An inverse problem cannot be posed
for a forward map that is not a function, so the scheduler must be carried explicitly
or held fixed and declared. We hold it fixed and declare it.

We note in passing that growth in S3 is trivially predictable: a right-hand side of $k$
edges gives growth of exactly $k$ per step ($1.00$, $2.00$, $3.00$, $4.00$ measured).
The interesting structure is not in the growth.

\section{Without intervention}\label{sec:passive}

Every identifiability win above came from \emph{setting the system's state}: choosing a
random initial row, choosing evaluation points. These are interventions in the sense of
Pearl~\cite{pearl2009}, not more careful passive observation, which is precisely why
they were cheap. Equation~\eqref{eq:decomp} carries an index for this reason, and this
section measures what the index costs.

We take the same $256$ automata on a $12$-cell ring and forbid the observer to touch
anything: nature picks the initial state, and the observer arrives---possibly long
after the transient has died---and watches. Computed exactly over all $4096$ starts for
every rule.

\begin{table}[h]
\centering
\begin{tabular}{lrr}
\toprule
Regime & Bits unresolved & Fully identified \\
\midrule
Intervene: choose the state, watch one step & 0.00 & 100\% \\
Passive: nature picks, observer arrives at $t=0$ & 0.51 & 66.9\% \\
Passive: nature picks, observer arrives late & 2.96 & 20.6\% \\
\bottomrule
\end{tabular}
\caption{The cost of not being allowed to intervene. A state handed to the observer and
watched for one step leaves $1.46$ bits; a state the observer chose leaves none.}
\end{table}

\paragraph{Why watching longer stops helping.}
A passive observer's entire future is one attractor period; past that the system
reshows states it has already shown. On this ring $30.2\%$ of starts land on a fixed
point and $42.0\%$ on a period of $1$ or $2$ (median period $6$, longest $240$). The
cap on identifiability is the \emph{diversity} of the attractor, not the duration of
the watch.

\paragraph{The worst case is total.}
Rules $238$--$255$ leave $7.00$ bits: $128$ rules permanently indistinguishable. Rule
$255$ drives every state to all-ones and remains there, so the neighbourhood $111$ is
the only one that ever occurs and seven of the eight table entries are never tested.

\paragraph{Removing an idealisation.}
Modelling ``nature gives $p$ points'' as a uniform random $p$-subset silently assumes
the observation points are spread out, whereas real attractors cluster. We therefore
replaced the model with an explicit driving map, so the observation points are whatever
its attractor visits.

\begin{table}[h]
\centering
\begin{tabular}{lrrr}
\toprule
Model of nature & Mean $p$ & Bits (late) & Identified \\
\midrule
Random map on the domain & 2.61 & 1.24 & 69\% \\
\quad \emph{same $p$, points spread uniformly} & 2.61 & 1.23 & 70\% \\
Affine contraction $\mathrm{round}(ax+b)$ & 1.23 & 2.56 & 39\% \\
\quad \emph{same $p$, points spread uniformly} & 1.23 & 2.67 & 40\% \\
\bottomrule
\end{tabular}
\caption{The clustering idealisation survives: spreading the points out is worth $0.01$
bits in one case and $-0.11$ in the other.}
\end{table}

The clustering hypothesis was wrong, and the null result is worth stating plainly:
conditioned on $p$, where the attractor sits barely matters. The damage is in $p$
itself. A dissipative system has a single fixed point, so $77\%$ of the time the late
observer sees one value and then that same value forever, and identifiability is
settled entirely by where it lands---$39\%$ of the time somewhere injective, and
otherwise nowhere useful, ever.

\paragraph{Consequence.} Passive identifiability is not a smaller number; it is a
distribution. With intervention, $p$ is a threshold and past it one always wins.
Without, $p$ only sets the odds, and $p$ is set by the system's own dynamics rather
than by the observer's patience. In a single-realisation science---one universe, one
economic history---the mean is not available either. One gets one draw.

\section{The generator-to-behaviour landscape has no local structure}\label{sec:landscape}

Recovering an elementary cellular automaton rule is trivial: from a random initial
condition the exact rule was recovered for all $256$ rules in a single linear pass over
the spacetime diagram. No search, no minimum-description-length machinery, no symbolic
regression.

And the recovered rule licenses no prediction. Taking compressed spacetime size as a
computable proxy for output complexity, rule $110$ scores $1562$ while its eight
single-bit neighbours score $145$, $146$, $151$, $308$, $329$, $440$, $783$ and $971$
--- between $0.09\times$ and $0.62\times$. Across all $2048$ single-bit rule mutations
the median change is $122$ against a full range of $3083$, and $10.7\%$ of mutations
move complexity by more than a quarter of that range.

The practical consequence is a negative one. A complexity landscape this rough cannot
be searched by gradient, hill-climbing, evolutionary, or any other neighbourhood
method, because proximity in rule space carries no information about proximity in
behaviour. Proposals for an ``automated primitive-search engine'' that hunts for
minimal generators must contend with the fact that, in the richest of our systems,
there is no landscape to climb---and separately, that the answer it would return is
free anyway.

\section{A fourth system, pre-registered}\label{sec:stochastic}

The three systems above are deterministic, fully observed, and algorithmic. If the
separation of \S\ref{sec:quadrants} is an artefact of that setting, a stochastic,
partially observed family should break it.

\paragraph{Exploratory run.}
A first family---all $3375$ transition matrices on three hidden states with rows drawn
from $\{0,\frac14,\frac12,\frac34,1\}$, lumped $\{0,1\}\!\to\!a$, $\{2\}\!\to\!b$---was
run against three criteria fixed in advance. Two passed; the quadrant-support criterion
\emph{failed}, because under passive observation the identifiable row was empty (one
generator of $3375$). Allowing the observer to reset the hidden state populated all four
quadrants, but that was a change of protocol made after seeing the failure, and we report
it only as the motivation for what follows.

\paragraph{Confirmatory run.}
We therefore froze a second design before running it: all $35^4 = 1{,}500{,}625$
transition matrices on \emph{four} hidden states, quarter-grid rows, a balanced $2+2$
lumping, horizon $T=6$, enumerated completely. Passive means a uniform initial state;
intervening means the experimenter may initialise each hidden state separately and pool
the four resulting laws. Metrics are the per-generator coordinates $i_O(g)$ and $p_1(g)$ of
\S\ref{sec:framework}, with the quadrant thresholds carried over unchanged
($i_O = 0$; $p_1 < 0.10$ bits/step). Five criteria were fixed: quadrant support at $1\%$
under intervention (C1); a $10$--$90$ percentile spread of $p_1$ of at least $0.50$ bits
among identifiable generators (C2); a spread of $i_O$ of at least $1$ bit among generators
at essentially fixed $p_1$ (C3); an exact finite-horizon observational-equivalence witness with
$p_1<0.10$ (C4), meaning generators whose $T=6$ observable laws coincide exactly; and at least $1\%$ of generators changing quadrant between protocols
(C5). C2 and C3 are distributional rather than existence tests, since over $1.5$M
generators an existence claim is nearly free.

\begin{table}[h]
\centering
\begin{tabular}{lrrrr}
\toprule
Protocol & id.\ \& pred. & id.\ \& not pred. & not id.\ \& pred. & not id.\ \& not pred. \\
\midrule
Passive      & $0.00\%$ & $0.00\%$ & $5.43\%$ & $94.57\%$ \\
Intervening  & $1.85\%$ & $80.16\%$ & $3.58\%$ & $14.41\%$ \\
\bottomrule
\end{tabular}
\caption{Quadrant occupancy over all $1{,}500{,}625$ generators. Passive $H(G\mid O)$ is
$3.24$ bits over $303{,}229$ distinct laws; intervening, $1.37$ bits over $1{,}235{,}809$.}
\end{table}

All five criteria were met. C2 gave a spread of $0.55$ bits, C3 a spread of $12.00$ bits,
C4 found $81{,}482$ generators sitting in non-singleton equivalence classes under the
exact $T=6$ observable law while satisfying $p_1 < 0.10$,
and C5 recorded $82.0\%$ quadrant migration.

Three things are worth drawing out. First, \emph{not one generator of $1{,}500{,}625$ is
uniquely identified by the passive observation map at horizon $T=6$}---the same empty row
that failed the exploratory criterion, reproduced on an unrelated family. We have not
shown that these generators are structurally non-identifiable under unlimited passive
data; the claim is about the declared observation map, and a longer horizon could only
shrink the classes. Second, the upper-right quadrant is
here populated for a reason entirely unlike the automaton case: the generator is known
exactly and the outcome is still uncertain because the law is stochastic, not because its
behaviour is undecidable. The same cell is reached by aleatoric and by computational
routes, which is evidence that the cell is a real category rather than an artefact of one
mechanism. Third, and we think most consequentially, C5 at $82.0\%$ says quadrant membership is
largely not a property of the system. The pipeline is
\[
  (S, \mathcal{A}, \Sigma, O)
  \;\longrightarrow\;
  \big(i_O(g),\, p_k(g)\big)
  \;\xrightarrow{\ \text{thresholds}\ }\;
  Q ,
\]
so varying $\mathcal{A}$ while holding $S$ fixed moves $Q$ for four generators in five.
The consequence for how the taxonomy should be read is direct: \emph{identifiable} and
\emph{predictable} are not intrinsic labels attached to a mechanism. They describe a
mechanism relative to an observation and intervention regime and a declared prediction
target. The intervention result of \S\ref{sec:passive} is the special case of this in
which only the number of bits was tracked, not the regime.

The empty passive row should not be read as a discovery. That state aggregation creates
observational equivalence classes, and that the inverse lumping problem has no unique
solution without further conditions, are established results. Ito, Amari and
Kobayashi~\cite{ito1992} solve precisely the question of when two different Markov chains
generate the same observed process, by linear algebra on the block structure of the
transition matrix. What the enumeration adds is the exhaustive
quantity---not one generator in $1{,}500{,}625$, over a completely enumerated class rather
than a constructed example---and the contrast with the interventional protocol on the same
family. The sharper question, under what conditions on the aggregation the identifiable
set is empty rather than merely small, we leave open and do not claim.

One exact witness from the exploratory family states the point most sharply. Five
distinct transition matrices induce an identical observable law \emph{at the horizon
observed}, with predictive entropy of exactly $0.000$ bits per step: $2.32$ bits of
uncertainty about the cause alongside zero bits of uncertainty about the consequence.
Whether the five remain equivalent at every horizon we did not test, and do not claim.

\section{Two distinct axes}\label{sec:quadrants}

Identifiability theory grew up in settings where recovering the parameters hands one
the trajectory, and treats recovery and prediction as a single achievement.

The claim of this paper rests on \S\ref{sec:stochastic} alone. There, a single completely
enumerated hypothesis class of $1{,}500{,}625$ generators populates all four combinations
of $i_O(g)$ and $p_1(g)$ under one fixed pair of definitions, with criteria frozen before
the run---$1.85\%$, $80.16\%$, $3.58\%$ and $14.41\%$ of the class in the four cells. That
is the evidence for \eqref{eq:nonimp}, and it does not depend on comparing systems to each
other.

The three deterministic systems are \emph{not} used to establish the entropy-based
quadrant claim, because for them that coordinate is degenerate: with the generator and the
current state both known, a deterministic system has $p_k \equiv 0$ by construction, so an
elementary cellular automaton is trivially ``predictable'' in the sense of $p_k$. What they
illustrate instead is that several distinct things get conflated under the word
prediction, and that each fails independently of generator recovery.

\begin{table}[h]
\centering
\begin{tabular}{p{2.9cm}p{3.0cm}p{7.6cm}}
\toprule
System & Sense of ``prediction'' & What generator recovery does not buy \\
\midrule
Bounded polynomials & exact future value & nothing: recovery does deliver prediction here.
This is the classical case. \\
\addlinespace
Elementary cellular automata & \emph{computational} predictability & the generator is
exactly recoverable in one linear pass, while the generator-to-complexity landscape is
highly non-smooth ($10.7\%$ of single-bit mutations move compressed output size by over a
quarter of its full range) and the rule family contains a computationally universal
member~\cite{cook2004}. Iterating a finite ring is of course always possible; what is not
available is any shortcut guided by proximity in rule space. \\
\addlinespace
Hypergraph rewriting & \emph{forward determinacy} & $1.09$ bits about which rule are
unresolved while edge growth is exactly $k$ per step; and with the scheduler free, the
outcome is not a function of the rule at all. \\
\addlinespace
Markov family (\S\ref{sec:stochastic}) & $p_1(g)$, entropy of the next observation &
the load-bearing case: all four cells occupied under one pair of measures. \\
\bottomrule
\end{tabular}
\caption{The deterministic systems illustrate distinct reasons why generator recovery,
behavioural computation and forward determinacy should not be conflated. They are not
evidence for the entropy-based quadrant claim, which rests on the last row.}
\end{table}

We state the resulting claim narrowly, and more narrowly than an earlier draft did.
Separating a structural axis from a randomness axis and plotting the pair over an
enumerated space of systems is not new: it is the complexity-entropy diagram of
computational mechanics~\cite{feldman2008}, discussed in \S\ref{sec:related}. The
tempting broad version of our claim---that no existing field has considered
identifiability and prediction separately---is an unsupportable negative. So is a claim of statistical independence: across the stochastic family the two
axes are weakly \emph{anti}-correlated ($r = -0.30$), which is precisely why we claim
distinctness rather than independence. What we demonstrate is this:
\begin{quote}
Within one finite-preimage framework and one pair of measures, exact enumeration
produces systems occupying all four combinations of the per-generator coordinates
$i_O(g)$ and $p_k(g)$; neither coordinate determines the other, in the sense of
\eqref{eq:nonimp}; and the choice of admissible observation class and of scheduler moves
a single system between those regimes.
\end{quote}
The off-diagonal cells are the interesting ones, and the reason they are rarely
confronted is structural rather than an oversight: the settings where identifiability
theory is usually applied are ones in which recovering the parameters does deliver the
trajectory, so the question of an exactly known and behaviourally useless generator does
not arise there. Since all four cells are occupied inside one hypothesis class under one
pair of measures, the thesis needs no cross-system comparison to stand, and we do not rest
it on one.

\section{Related work}\label{sec:related}

The framing of \S\ref{sec:framework} is structural identifiability
\cite{bellman1970,walter1997}; the equivalence-class formulation is standard in causal
discovery, where algorithms return a Markov equivalence class rather than a graph
\cite{verma1990,pearl2009}, and in automata theory via Myhill--Nerode
\cite{nerode1958}. Optimal experimental design by information gain is due to Lindley
\cite{lindley1956}, with the modern active-learning form in \cite{houlsby2011}.
Generative reachability under a primitive set is universal algebra: Post's lattice
\cite{post1941} classifies the Boolean case completely. Integer complexity is
\cite{oeis005245} and \cite[\S F26]{guy2004}. Universality of rule 110 is due to Cook
\cite{cook2004}; the broader automaton programme is \cite{wolfram2002}. Causal
invariance in rewriting systems is \cite{wolfram2020}. The Weisfeiler--Leman invariant
is \cite{weisfeiler1968}, its incompleteness \cite{cai1992}. Algorithmic information
theory, which sets the ceiling on compressibility arguments generally, is
\cite{livitanyi2008}.

The closest prior work, and the comparison a reader should make first, is computational
mechanics. Feldman, McTague and Crutchfield~\cite{feldman2008} plot a measure of
structural complexity against a measure of randomness over an enumerated space of
processes---maps of the interval, cellular automata, Ising systems, Markov chains---which
is the same two-axis move, made eighteen years earlier and over a comparably broad
sweep of systems. Our vertical coordinate $p_k$ is, in the stationary limit, their
entropy rate $h_\mu$.

The horizontal coordinates differ, and the difference is the whole of what we add. Their
statistical complexity $C_\mu$ is the memory of the minimal predictor and is a property of
the \emph{process}; the $\epsilon$-machine is precisely the canonical representative of an
observational equivalence class, and computational mechanics proceeds by quotienting that
class away. Our $i_O(g)$ measures the class itself---how many members of a declared
hypothesis space $\HH$ remain indistinguishable under a declared observation map $O$---and
is therefore a property of $g$ relative to $(\HH, O)$ rather than of the process. A process
can have small $C_\mu$ and large $i_O$, which is exactly the lower-left quadrant of
\S\ref{sec:quadrants}: a simple minimal predictor that many distinct mechanisms share.

That interventions change identifiability is itself well established: active learning of
causal models is built on choosing interventions that guarantee it, with optimal strategies
known~\cite{hauser2014}. Our claim is correspondingly narrow. We are not aware of prior
work that combines per-generator observational-equivalence size with predictive entropy and
exhaustively measures how the resulting joint classification migrates under a change in
admissible interventions---$Q = Q(S, \mathcal{A}, \Sigma, O)$, at $82.0\%$ in
\S\ref{sec:stochastic}. We make no stronger claim than that.

\section{Limitations}

Identifiability in \S\ref{sec:stochastic} is equality of the finite-dimensional law at
horizon $T=6$, not of the full process law. Every preimage reported there is therefore an
upper bound on the preimage available to an observer with unlimited passive data, and the
empty passive row is a statement about that observation map rather than about structural
identifiability in the sense of~\cite{ito1992}.
The stochastic results of \S\ref{sec:stochastic} separate cleanly into an exploratory
run, whose interventional analysis was post-hoc and is reported as such, and a
confirmatory run frozen before execution. Only the latter supports the claim. The
quadrant thresholds ($i_O(g)=0$, $p_1(g)<0.10$) are conventions carried between the two
rather than optimised, but they remain conventions.
An earlier version of this manuscript reported $\log_2\mathbb{E}[|P(O)|]$ for two of
the three systems and $H(G\mid O)$ for the third; the figures here are $H(G\mid O)$
throughout, which lowers every affected number (the rewriting residual from $1.55$ to
$1.09$ bits) without changing any qualitative conclusion.
The hypothesis spaces are small and finite by design, which is what permits exact
enumeration; nothing here is an asymptotic claim. S3 comparisons use a
Weisfeiler--Leman invariant, so the $1.09$-bit residual is a lower bound on
distinguishability and could only grow under a complete test. S3 results hold the
scheduler seed fixed, which controls the order-dependence of \S\ref{sec:scheduler}
rather than removing it. The $26{,}624$ figure for minimal expressions of $107$ counts
ordered trees and so overcounts by commutativity, though not by an order of magnitude.
Counts of permanently unidentifiable automata are mildly ring-width sensitive ($132$ at
width $64$, $134$ at width $201$) and should always be reported with the width. The
compressed-size proxy for output complexity in \S\ref{sec:landscape} is a computable
stand-in for an uncomputable quantity and is used only for relative comparison between
neighbouring rules.

\section{Discussion}

The programme this paper argues against---search for the smallest rule generating an
observed object---fails in two independent ways. In the simplest system it is
ill-posed: minima are non-unique, unstable, and lack a closed form. In the richest it
is trivial and unhelpful: the rule is free and predicts nothing.

What replaces it is narrower and testable. Decompose $H(G\mid O)$ by
Equation~\eqref{eq:decomp}, with the admissible class $\mathcal{A}$ stated explicitly
rather than left implicit, and report which quadrant of \S\ref{sec:quadrants} the system occupies. All
three are measurable. The first of those questions has now been answered
once: \S\ref{sec:stochastic} carries the classification into stochastic, partially
observed dynamics with a hypothesis class four orders of magnitude larger, under criteria
fixed in advance. What remains untested is continuous state and continuous time, where
exact enumeration is unavailable and the preimage must be estimated rather than counted.
We would want that before claiming the quadrant structure is general.

The broader statement the results now support is that a system has no single degree of
knowability. What can be known depends on which question is asked---which mechanism is
acting, or what it will do next---and on what the observer is permitted to do to it.

\paragraph{Reproducibility.}
All figures are produced by seven self-contained Python scripts requiring only the
standard library: \texttt{mgs.py} (algebra and reachability), \texttt{mgs\_ca.py}
(cellular automata), \texttt{mgs\_rewrite.py} (hypergraph rewriting),
\texttt{mgs\_ident.py} (identifiability curves), \texttt{mgs\_passive.py} (passive
observation), \texttt{mgs\_stochastic.py} (the exploratory Markov family) and
\texttt{mgs\_preregistered.py} (the frozen confirmatory run, which additionally requires
NumPy). All but the last run in a few minutes on one core.

No number in this manuscript is manually trusted. \texttt{audit.py} re-derives every
quantitative assertion from the script that produces it---recomputing from scratch the few
that no committed script emitted---and reports each against the value printed here. It
tags provenance as exhaustive computation, sampled computation, derivation, or recomputed
in the audit itself, so a reader can see which claims rest on enumeration and which on a
fixed seed. The current run is $45$ of $45$ claims verified, recorded in
\texttt{AUDIT.txt}, with the values also emitted to \texttt{RESULTS.txt} in a form
\LaTeX{} can import directly. Bibliographic and interpretive assertions are outside its
scope and were checked by hand.

\begin{thebibliography}{99}
\bibitem{bellman1970} R. Bellman and K. J. {\AA}str\"om. On structural identifiability.
\emph{Mathematical Biosciences}, 7(3--4):329--339, 1970.
\bibitem{cai1992} J.-Y. Cai, M. F\"urer, and N. Immerman. An optimal lower bound on the
number of variables for graph identification. \emph{Combinatorica}, 12(4):389--410, 1992.
\bibitem{feldman2008} D. P. Feldman, C. S. McTague, and J. P. Crutchfield. The
organization of intrinsic computation: complexity-entropy diagrams and the diversity of
natural information processing. \emph{Chaos}, 18(4):043106, 2008.
\bibitem{cook2004} M. Cook. Universality in elementary cellular automata.
\emph{Complex Systems}, 15(1):1--40, 2004.
\bibitem{hauser2014} A. Hauser and P. B\"uhlmann. Two optimal strategies for active
learning of causal models from interventional data. \emph{International Journal of
Approximate Reasoning}, 55(4):926--939, 2014.
\bibitem{ito1992} H. Ito, S.-I. Amari, and K. Kobayashi. Identifiability of hidden Markov
information sources and their minimum degrees of freedom. \emph{IEEE Transactions on
Information Theory}, 38(2):324--333, 1992.
\bibitem{jacquemard2003} F. Jacquemard. Reachability and confluence are undecidable for
flat term rewriting systems. \emph{Information Processing Letters}, 87(5):265--270, 2003.
\bibitem{guy2004} R. K. Guy. \emph{Unsolved Problems in Number Theory}. Springer, 3rd
edition, 2004.
\bibitem{houlsby2011} N. Houlsby, F. Husz\'ar, Z. Ghahramani, and M. Lengyel. Bayesian
active learning for classification and preference learning. arXiv:1112.5745, 2011.
\bibitem{lindley1956} D. V. Lindley. On a measure of the information provided by an
experiment. \emph{Annals of Mathematical Statistics}, 27(4):986--1005, 1956.
\bibitem{livitanyi2008} M. Li and P. Vit\'anyi. \emph{An Introduction to Kolmogorov
Complexity and Its Applications}. Springer, 3rd edition, 2008.
\bibitem{nerode1958} A. Nerode. Linear automaton transformations.
\emph{Proceedings of the American Mathematical Society}, 9(4):541--544, 1958.
\bibitem{oeis005245} OEIS Foundation. Sequence A005245: complexity of $n$, the number
of 1's required to build $n$ using $+$ and $*$. \url{https://oeis.org/A005245}.
\bibitem{pearl2009} J. Pearl. \emph{Causality: Models, Reasoning, and Inference}.
Cambridge University Press, 2nd edition, 2009.
\bibitem{post1941} E. L. Post. \emph{The Two-Valued Iterative Systems of Mathematical
Logic}. Annals of Mathematics Studies 5, Princeton University Press, 1941.
\bibitem{verma1990} T. Verma and J. Pearl. Equivalence and synthesis of causal models.
In \emph{Proceedings of the Sixth Conference on Uncertainty in Artificial
Intelligence}, pages 220--227, 1990.
\bibitem{walter1997} E. Walter and L. Pronzato. \emph{Identification of Parametric
Models from Experimental Data}. Springer, 1997.
\bibitem{weisfeiler1968} B. Weisfeiler and A. Leman. A reduction of a graph to a
canonical form and an algebra arising during this reduction.
\emph{Nauchno-Technicheskaya Informatsia}, 2(9):12--16, 1968.
\bibitem{wolfram2002} S. Wolfram. \emph{A New Kind of Science}. Wolfram Media, 2002.
\bibitem{wolfram2020} S. Wolfram. A class of models with the potential to represent
fundamental physics. \emph{Complex Systems}, 29(2):107--536, 2020. DOI 10.25088/ComplexSystems.29.2.107.
\end{thebibliography}

\end{document}
